\section{Hardness and Approximability}
\label{sec:hardness}

\subsection{NP-hardness of \CBSLAO-\textsc{dec}}

\begin{theorem}[NP-hardness, exact]
\label{thm:hardness}
\CBSLAO-\textsc{dec} is NP-hard, even when
(i) all cost distributions $\mathcal{D}_i^c$ are degenerate point
masses,
(ii) latency is unconstrained ($T = \infty$),
(iii) the utility function is additive across tool invocations,
and (iv) the policy class is $\Pi_{\mathrm{non}}$.
\end{theorem}

\begin{proof}
We reduce from the 0/1 Knapsack decision problem, which is NP-hard
\cite{garey1979computers}.

An instance of 0/1 Knapsack is
$(w_1, \dots, w_n, v_1, \dots, v_n, W, V)$ with
$w_i, v_i, W, V \in \mathbb{Z}_{\ge 0}$; we ask whether there exists
$S \subseteq [n]$ with
$\sum_{i \in S} w_i \le W$ and $\sum_{i \in S} v_i \ge V$.

Construct a \CBSLAO~instance as follows. Let
$\mathcal{S} = \{s_1, \dots, s_n\}$ with degenerate cost distributions
$\mathcal{D}_i^c = \delta_{w_i}$ (point mass at~$w_i$). Let latency be
zero and $T = \infty$. Define the utility as
\[
  U(h) \;=\; \min\!\Bigl\{1,\; \tfrac{1}{V}\sum_{i : s_i \in h} v_i\Bigr\}.
\]
Set $B = W$, $\delta = 0$, and $u^\star = 1$.

Since costs are deterministic, the chance constraint reduces to a hard
constraint $\sum_{i \in S} w_i \le W$. Since utility is
saturating-additive and $u^\star = 1$, a policy achieves utility
$\ge u^\star$ iff the chosen subset satisfies
$\sum_{i \in S} v_i \ge V$. Any non-adaptive policy chooses a subset
$S \subseteq [n]$. Therefore a \CBSLAO~policy achieving $u^\star = 1$
within budget $B = W$ exists iff a knapsack solution of value $\ge V$
within weight $\le W$ exists. The reduction is polynomial in input
size.
\qed
\end{proof}

\begin{remark}
\label{rem:har}
Because the reduction uses only $\Pi_{\mathrm{non}}$, we have
\emph{a fortiori} hardness for $\Pi_{\mathrm{semi}}$ and
$\Pi_{\mathrm{ad}}$. Moreover, the reduction is purely
deterministic, so \CBSLAO-\textsc{dec} is NP-hard already in the
deterministic regime --- stochasticity in real deployments can only
make the problem harder. Theorem~\ref{thm:hardness} rules out exact
polynomial-time solution but, since 0/1 Knapsack itself admits an
FPTAS, leaves open the possibility of a polynomial-time approximation.
The next subsection closes that loophole.
\end{remark}

\subsection{Hardness of approximation}
\label{sec:inapprox}

The 0/1 Knapsack reduction of Theorem~\ref{thm:hardness} only certifies
NP-hardness of \emph{exact} solution. We now show a stronger
\emph{inapproximability} result by reducing from
Maximum~$k$-Coverage~\cite{feige1998threshold}, exploiting the
schema-typing structure of \CBSLAO~tools rather than the
additive-utility structure of Knapsack.

\begin{theorem}[Inapproximability of \CBSLAO-\textsc{opt}]
\label{thm:inapprox}
Unless $\mathrm{P} = \mathrm{NP}$, no polynomial-time algorithm
achieves a $(1 - 1/e + \varepsilon)$-approximation to \CBSLAO-\textsc{opt}
for any fixed $\varepsilon > 0$, even when
(i) costs are deterministic,
(ii) latency is unconstrained,
(iii) the policy class is $\Pi_{\mathrm{non}}$, and
(iv) the utility function is monotone submodular.
\end{theorem}

\begin{proof}
We reduce from Max-$k$-Coverage, which is NP-hard to approximate within
$(1 - 1/e + \varepsilon)$ unless $\mathrm{P} = \mathrm{NP}$ by
Feige's threshold~\cite{feige1998threshold}.

An instance of Max-$k$-Coverage is a triple
$(\mathcal{U}, \mathcal{F}, k)$ with universe $\mathcal{U}$, family of
sets $\mathcal{F} = \{S_1, \dots, S_m\}$ where $S_j \subseteq \mathcal{U}$,
and budget $k \in \mathbb{N}$; the goal is to maximise
$|\bigcup_{j \in A} S_j|$ subject to $|A| \le k$.

Construct a \CBSLAO~instance as follows. Let the type universe be
$\mathcal{T} = \{\tau_u : u \in \mathcal{U}\} \cup \{\tau_0\}$. For each
set $S_j \in \mathcal{F}$ create a tool $s_j$ with
$\mathsf{In}_j = \{\tau_0\}$, $\mathsf{Out}_j = \{\tau_u : u \in S_j\}$,
deterministic cost~$1$, latency~$0$. The query $q$ has type
$\{\tau_0\}$, $B = k$, $T = \infty$, $\delta = 0$. Define the utility
\[
  U(h) \;=\; \frac{1}{|\mathcal{U}|}
  \Bigl|\!\bigcup_{s_j \in h} S_j\Bigr|.
\]
The set function $A \mapsto |\bigcup_{j \in A} S_j|$ is monotone
submodular~\cite{nemhauser1978best}; therefore $U$ is monotone
submodular. The chance constraint with $\delta = 0$ becomes
$|A| \le k$.

Hence the optimal non-adaptive \CBSLAO~value equals
$|\mathcal{U}|^{-1}$ times the optimal Max-$k$-Coverage value. A
polynomial-time $(1 - 1/e + \varepsilon)$-approximation for the
\CBSLAO~instance therefore yields a
$(1 - 1/e + \varepsilon)$-approximation for Max-$k$-Coverage,
contradicting~\cite{feige1998threshold}. The reduction is polynomial
in the input size.
\qed
\end{proof}

\begin{remark}
\label{rem:inapprox-tightness}
Theorem~\ref{thm:inapprox} subsumes the trivial inapproximability of
Corollary~\ref{cor:inap-trivial} below: any polynomial-time algorithm
that violated the $(1 - 1/e + \varepsilon)$-approximation bound would,
\emph{a fortiori}, also achieve exact optimality on the constructed
instance. The reduction further uses the type-closure structure of
schemas: it is the typing of $\mathsf{In}_j, \mathsf{Out}_j$ that lets
the reduction encode element-set incidence as a service-composition
constraint, distinguishing the result from generic
stochastic-knapsack hardness.
\end{remark}

\begin{corollary}[Trivial-exact inapproximability, for completeness]
\label{cor:inap-trivial}
Unless $\mathrm{P} = \mathrm{NP}$, no polynomial-time algorithm returns
a feasible policy whose value matches the optimum utility exactly for
every \CBSLAO~instance.
\end{corollary}

\subsection{A matching $(1-1/e)$-approximation in the offline non-adaptive setting}
\label{sec:approx}

The bound of Theorem~\ref{thm:inapprox} is \emph{tight}: the offline
non-adaptive variant of \CBSLAO~admits a polynomial-time
$(1-1/e)$-approximation when the utility function is monotone
submodular, matching the inapproximability lower bound up to the
$\varepsilon$ slack. This justifies the \emph{greedy-offline} phase of
\CBUC~with a formal guarantee rather than a heuristic.

\begin{theorem}[Offline approximation]
\label{thm:approx}
For the offline non-adaptive version of \CBSLAO~with deterministic
costs and saturating-additive utility
$U(h) = g\!\bigl(\sum_{i \in A(h)} u_i\bigr)$, where
$g: \mathbb{R}_{\ge 0} \to [0, 1]$ is monotone non-decreasing and
concave with $g(0) = 0$, the modified-greedy algorithm of
Sviridenko~\cite{sviridenko2004} returns a feasible policy whose value
is at least $(1 - 1/e) \cdot \mathsf{OPT}$ in time $\Oh(K^4)$.
\end{theorem}

\begin{proof}
The set function $f(A) = g\!\bigl(\sum_{i \in A} u_i\bigr)$ is monotone
since $g$ is monotone and $u_i \ge 0$. It is submodular because the
composition of a non-decreasing concave function $g$ with a
non-negative additive set function $A \mapsto \sum_{i \in A} u_i$ is
monotone submodular: for $A \subseteq B$ and $i \notin B$,
\[
  f(A \cup \{i\}) - f(A)
   = g(s_A + u_i) - g(s_A)
   \;\ge\; g(s_B + u_i) - g(s_B)
   = f(B \cup \{i\}) - f(B),
\]
where $s_A = \sum_{j \in A} u_j$ and the inequality follows from the
concavity of $g$ (cf.~\cite[Prop.~2.2.6]{nemhauser1978best}).

Maximising a monotone submodular function under a knapsack constraint
is exactly Sviridenko's setting~\cite{sviridenko2004}, which gives the
$(1-1/e)$-approximation in $\Oh(K^4)$ time via a partial-enumeration
modified greedy.
\qed
\end{proof}

\begin{remark}
The simulator's saturating utility $U(h) = 1 - \exp(-\sum_i u_i)$
takes $g(x) = 1 - e^{-x}$, which is monotone, concave, and satisfies
$g(0) = 0$, so Theorem~\ref{thm:approx} applies. Together,
Theorems~\ref{thm:inapprox} and~\ref{thm:approx} pin down the
polynomial-time approximability of the offline non-adaptive variant
exactly at $1 - 1/e$.
\end{remark}

\subsection{Adaptivity gap}

The regret bound of \S\ref{sec:algorithm} (Theorem~\ref{thm:regret})
targets the optimal non-adaptive feasible policy
$\pi^\dagger_{\mathrm{CC}}$. We now show that the gap between
$\pi^\dagger_{\mathrm{CC}}$ and the adaptive optimum is bounded above
by a constant on every instance, but is strictly larger than~$1$ on
some instances.

\begin{lemma}[Adaptivity gap, explicit construction and constant ceiling]
\label{lem:adaptivity-gap}
There exists a family of \CBSLAO~instances $\{I_n\}_{n \ge 2}$ on which
\[
  \mathsf{OPT}_{\mathrm{ad}}(I_n)
   \ \ge\ \tfrac{4}{3} \cdot
       \mathsf{OPT}_{\mathrm{non}}(I_n) \cdot (1 - o(1)).
\]
Conversely, on every \CBSLAO~instance with deterministic latency and
additive utility, $\mathsf{OPT}_{\mathrm{ad}}(I) \le c \cdot
\mathsf{OPT}_{\mathrm{non}}(I)$ for an absolute constant
$c \le 4$~\cite{dean2008stochastic}.
\end{lemma}

\begin{proof}[sketch]
\textbf{Lower bound (gap is at least $4/3$).}
Embed the $4/3$-construction of Dean, Goemans, and
Vondr\'{a}k~\cite[Sec.~5]{dean2008stochastic} for stochastic knapsack
into \CBSLAO~by setting $\delta = 0$, latency to zero, and utility to
the additive Knapsack value. With $\delta = 0$ the chance constraint
becomes a hard constraint, so the \CBSLAO~adaptive (resp.\
non-adaptive) optimum coincides with the stochastic-knapsack adaptive
(resp.\ non-adaptive) optimum on the same instance. The gap factor of
$4/3$ as $n \to \infty$ transfers verbatim. Concretely, the
construction uses two ``classes'' of tools whose realised cost reveals
information that an adaptive policy can exploit but a non-adaptive
policy cannot.

\textbf{Upper bound (constant ceiling).}
\cite[Thm.~1]{dean2008stochastic} establishes a constant adaptivity
gap of at most $4$ for one-resource stochastic knapsack with additive
value. The deadline chance-constraint of \CBSLAO~adds a second
resource of identical type; intersecting two stochastic-knapsack
constraints multiplies the gap by at most another constant via a
routine product argument (see also~\cite{dean2008stochastic}, Remark
following Thm.~1).
\qed
\end{proof}

\begin{remark}[Empirical adaptivity gap on the simulator]
\label{rem:gap-empirical}
Lemma~\ref{lem:adaptivity-gap} bounds the gap by an absolute constant
in the worst case. On our simulator the empirical ratio of
non-adaptive to informed-knapsack utility is
$1.00 \pm 0.05$ across all sweep cells (median $1.00$;
Remark~\ref{rem:adaptivegap}), well below the $4/3$ worst-case factor.
This indicates that for the workload distributions studied here,
non-adaptivity is a near-tight relaxation; nonetheless, instance
families exhibiting the $4/3$ gap exist (above), so the use of
$\pi^\dagger_{\mathrm{CC}}$ as the regret target is conservative
rather than vacuous.
\end{remark}
